% Intended LaTeX compiler: pdflatex
\documentclass[a4paper,12pt]{article}
\usepackage[utf8]{inputenc}
\usepackage[T1]{fontenc}
\usepackage{graphicx}
\usepackage{longtable}
\usepackage{wrapfig}
\usepackage{rotating}
\usepackage[normalem]{ulem}
\usepackage{amsmath}
\usepackage{amssymb}
\usepackage{capt-of}
\usepackage{hyperref}
\usepackage{\string~/.emacs.d/common}
\author{mahmood}
\date{\today}
\title{graph}
\hypersetup{
 pdfauthor={mahmood},
 pdftitle={graph},
 pdfkeywords={},
 pdfsubject={},
 pdfcreator={Emacs 28.2 (Org mode 9.5.5)}, 
 pdflang={English}}
\begin{document}

\maketitle
\begin{definition}
a \textbf{graph} is an \href{discrete_maths2.org}{ordered pair} \(G = (V, E)\), where \(V\) is a set whose elements are called \href{20230314001051-graph.org}{vertices}, and \(E\) is a set of paired vertices, whose elements are called \href{20230314001632-edge.org}{edge}s\\
\end{definition}
\begin{my_example}
consider the graph \(G_1=(V_1,E_1),V_1=\{a,b,c\},E_1=\{(a,b),(a,c),(a,a),e_1=(b,c),e_2=(b,c)\)\\
\includesvg{/home/mahmooz/brain/out/Eqzl9lY}\\
\end{my_example}
\begin{my_example}
consider the following \href{20220706211958-algorithm.org}{algorithm} which receives a non-directed graph and returns a list of its edges\\
\includesvg{/home/mahmooz/brain/out/1IQAdwQ}\\

the \href{20221130014441-time_complexity.org}{time complexity} of this algorithm is \(O(|E|+|V|)\)\\
\end{my_example}
\begin{lemma}
the maximum number of edges in a simple non-directed graph is\\
\[
\frac{|V| \cdot (|V|-1)}{2}
\]\\
\begin{proof}
in a complete graph, the rank of each vertex is \(|V|-1\), and it follows trivially that the sum of ranks is \(|V|(|V|-1)\), and since each vertex touches 2 edges, the number of edges is \(|E|=\frac{|V| \cdot (|V|-1)}{2}\)\\
\end{proof}
\end{lemma}
\begin{definition}
let \(G=(V,E)\) be a graph, a \textbf{path} from the a \(v_0\) to another \(v_1\) is a sequence of edges\\
\[
\alpha = (v_0,v_1),(v_1,v_2),\dots,(v_{l-1},v_l)
\]\\
such that the first edge is between \(v_0\) and \(v_1\), the second is between \(v_1\) and \(v_2\), and in general, for all \(i\), \(1 \le i \le l\)\\
the \(i\)'th edge is between \(v_{i-1}\) and \(v_i\). we can also write \(a=(v_0,v_1,\dots,v_l)\)\\
a path in a \href{20230314002402-directed_graph.org}{directed graph} is the same as a graph of a non-directed graph, with one additional constraint, all the edges should be directed in the same direction, from the first to the last one\\
\end{definition}
\begin{definition}
let \(\alpha=(v_0,v_1,\dots,v_l)\) be a path in a graph \(G\)\\
\begin{enumerate}
\item if \(v_0=v_l\), then \(a\) is a \textbf{cycle}\\
\item a path \(\alpha\) is \textbf{simple} if no vertex appears in \(\alpha\) more than once\\
\item a cycle \(\sigma\) is simple if no vertex except for the first one appears in \(sigma\) more than once\\
\item let \(\alpha=(v_1,v_2,\dots,v_l),\beta=(v_l,v_{l+1},\dots,v_{l+m})\) be two paths in the graph \(G=(V,E)\), the concatenation of \(\alpha,\beta\) is \(\alpha \circ \beta = (v_0,v_1,\dots,v_l,v_{l+1},\dots,v_{l+m})\)\\
\end{enumerate}
\end{definition}
\begin{definition}
\begin{enumerate}
\item the graph \(G=(V,E)\) is \textbf{connected} if between every two vertices \(u,v \in V\) there's a path\\
\item a graph \(G=(V,E)\) is \textbf{acyclic} if in \(G\) there's no cycles\\
\item a graph \(T=(V,E)\) is a tree if it is connected and acyclic\\
\end{enumerate}
\end{definition}
\begin{definition}
let \(G=(V,E)\) be a connected non-weighted graph and let \(\alpha=v_0,v_1,\dots,v_l\) be a path in \(G\)\\
\begin{enumerate}
\item the length of \(\alpha\) is defined as the number of edges in the path, \(l\), we write \(|\alpha|=l\)\\
\item the distance between the vertices \(u\) and \(v\) is defined as the length of the minimal path between \(u\) and \(v\), we denote by \(\delta_G(u,v)\) the distance in graph G between \(u\) and \(v\) and get:\\
\end{enumerate}
\[
\delta(u,v) = \min_a |a|
\]\\
if there is no path between \(u\) and \(v\) then \(\delta(u,v)=\infty\)\\
\end{definition}

\begin{my_example}
here we demonstrate an algorithm that finds the distances from a given vertex to all the other vertices in a non-weighted graph\\
we use the \textbf{breadth first search} technique along with a \href{data_structures.org}{queue}\\
\includesvg{/home/mahmooz/brain/out/Bbwoajo}\\
\end{my_example}
\end{document}